\documentclass[UTF8,12pt]{ctexart}
\usepackage[a4paper,margin=24mm]{geometry}
\usepackage{amsmath,amssymb,bm,graphicx,adjustbox}
\newcommand{\fitmath}[1]{\begin{adjustbox}{max width=\linewidth}$\displaystyle #1$\end{adjustbox}}
\setlength{\parindent}{0pt}
\setlength{\parskip}{.7em}
\sloppy
\allowdisplaybreaks
\begin{document}
\section*{2003 年考研数学三 · 参考答案}
\subsection*{2003年真题参考答案}

\subsection*{一、填空题}

（1）\(\displaystyle \lambda>2\)。（2）\(\displaystyle 4a^6\)。（3）\(\displaystyle a^2\)。（4）\(\displaystyle -1\)。（5）\(\displaystyle 0.9\)。（6）\(\displaystyle \dfrac{\displaystyle 1}{\displaystyle 2}\)。


\subsection*{二、选择题}

（1）D。（2）A。（3）B。（4）C。（5）B。（6）C。


\subsection*{三、}

令 \(\displaystyle f(1)=\dfrac{\displaystyle 1}{\displaystyle \pi}\)，则 \(\displaystyle f(x)\) 在 \(\displaystyle [\dfrac{\displaystyle 1}{\displaystyle 2},\allowbreak 1]\) 上连续。


\subsection*{四、}

\(\displaystyle \dfrac{\displaystyle \partial^2g}{\displaystyle \partial x^2}+\dfrac{\displaystyle \partial^2g}{\displaystyle \partial y^2}=x^2+y^2\)。


\subsection*{五、}

\(\displaystyle I=\dfrac{\displaystyle \pi}{\displaystyle 2}(1+e^\pi)\)。


\subsection*{六、}

\(\displaystyle f(x)=1-\dfrac{\displaystyle 1}{\displaystyle 2}\ln(1+x^2)\;(|x|<1)\)，\(\displaystyle f(x)\) 在 \(\displaystyle x=0\) 处取得极大值，极大值为 \(\displaystyle f(0)=1\)。


\subsection*{七、}

（1）\(\displaystyle F^{\prime}(x)+2F(x)=4e^{2x}\)。（2）\(\displaystyle F(x)=e^{2x}-e^{-2x}\)。


\subsection*{八、}

证明略。（可利用介值定理和罗尔定理。）


\subsection*{九、}

（1）当 \(\displaystyle b\ne0\) 且 \(\displaystyle b+\displaystyle \sum_{i=1}^na_i\ne0\) 时，方程组仅有零解。（2）当 \(\displaystyle b=0\) 时，方程组有非零解，它的一个基础解系为 \(\displaystyle \alpha_1=(-\dfrac{\displaystyle a_2}{\displaystyle a_1},\allowbreak 1,\allowbreak 0,\allowbreak \ldots,\allowbreak 0)^{\mathrm T}\)，\(\displaystyle \alpha_2=(-\dfrac{\displaystyle a_3}{\displaystyle a_1},\allowbreak 0,\allowbreak 1,\allowbreak \ldots,\allowbreak 0)^{\mathrm T}\)，\(\displaystyle \ldots\)，\(\displaystyle \alpha_{n-1}=(-\dfrac{\displaystyle a_n}{\displaystyle a_1},\allowbreak 0,\allowbreak 0,\allowbreak \ldots,\allowbreak 1)^{\mathrm T}\)。当 \(\displaystyle b=-\displaystyle \sum_{i=1}^na_i\) 时，方程组有非零解，它的一个基础解系为 \(\displaystyle (1,\allowbreak 1,\allowbreak \ldots,\allowbreak 1)^{\mathrm T}\)。


\subsection*{十、}

（1）\(\displaystyle a=1,\allowbreak b=2\)。（2）令 \(\displaystyle Q=\begin{pmatrix}\dfrac{\displaystyle 2}{\displaystyle \sqrt5}&0&\dfrac{\displaystyle 1}{\displaystyle \sqrt5}\\0&1&0\\\dfrac{\displaystyle 1}{\displaystyle \sqrt5}&0&-\dfrac{\displaystyle 2}{\displaystyle \sqrt5}\end{pmatrix}\)，\(\displaystyle Q\) 为正交矩阵。在正交变换 \(\displaystyle X=QY\) 下，二次型 \(\displaystyle f\) 化为标准形 \(\displaystyle f=2y_1^2+2y_2^2-3y_3^2\)。


\subsection*{十一、}

\(\displaystyle Y=F(X)\) 的分布函数 \(\displaystyle G(y)=\begin{cases}0,\allowbreak &y\le0,\allowbreak \\y,\allowbreak &0<y<1,\allowbreak \\1,\allowbreak &y\ge1.\end{cases}\)


\subsection*{十二、}

\(\displaystyle g(u)=0.3f(u-1)+0.7f(u-2)\)。

\end{document}
